%==============================================================================
%  CT-ATPT: Adaptive Token Pruning Transformer for Lung Nodule Malignancy
%  Classification on 3D CT
%
%  Target venue: IEEE Transactions on Medical Imaging / IEEE JBHI / IEEE Access
%  Class file  : IEEEtran (https://www.ctan.org/pkg/ieeetran)
%
%  Compile with:
%      pdflatex paper_ctatpt
%      bibtex   paper_ctatpt
%      pdflatex paper_ctatpt
%      pdflatex paper_ctatpt
%==============================================================================
\documentclass[journal]{IEEEtran}

%------------------------------------------------------------------------------
%  Packages
%------------------------------------------------------------------------------
\usepackage{cite}
\usepackage{amsmath,amssymb,amsfonts}
\usepackage{algorithmic}
\usepackage{graphicx}
\usepackage{textcomp}
\usepackage{xcolor}
\usepackage{booktabs}
\usepackage{array}
\usepackage{multirow}
\usepackage{url}
\usepackage{hyperref}
\hypersetup{
    colorlinks=true,
    linkcolor=blue,
    citecolor=blue,
    urlcolor=blue
}

%------------------------------------------------------------------------------
%  Result-placeholder macro.
%  Use \R{label} for any number that will be filled in once training finishes.
%  Example: \R{AUC} renders as a coloured [AUC] block so missing values are
%  immediately visible during proof-reading.
%------------------------------------------------------------------------------
\newcommand{\R}[1]{\textcolor{red}{\textbf{[#1]}}}
\newcommand{\todo}[1]{\textcolor{orange}{\textbf{TODO: #1}}}

%------------------------------------------------------------------------------
%  Math shortcuts
%------------------------------------------------------------------------------
\DeclareMathOperator*{\softmax}{softmax}
\DeclareMathOperator*{\Var}{Var}
\newcommand{\norm}[1]{\left\lVert#1\right\rVert}
\newcommand{\Lcls}{\mathcal{L}_{\mathrm{cls}}}
\newcommand{\Lent}{\mathcal{L}_{\mathrm{ent}}}
\newcommand{\Lcon}{\mathcal{L}_{\mathrm{con}}}
\newcommand{\Lspar}{\mathcal{L}_{\mathrm{spar}}}

%==============================================================================
\begin{document}

\title{CT-ATPT: An Adaptive Token Pruning Transformer for Lung Nodule
Malignancy Classification on 3D CT}

\author{
    \IEEEauthorblockN{%
        \R{Author~1}, \R{Author~2}, \IEEEmembership{Student~Member,~IEEE},
        and \R{Author~3}, \IEEEmembership{Member,~IEEE}}
    \thanks{\R{Author~1} and \R{Author~2} are with the
        \R{Department of \ldots, University of \ldots, City, Country}
        (e-mail: \R{a1@example.edu; a2@example.edu}).}%
    \thanks{\R{Author~3} is with \R{Department, University, City, Country}
        (e-mail: \R{a3@example.edu}).}%
    \thanks{Manuscript received \R{Month~DD, YYYY}; revised
        \R{Month~DD, YYYY}.}%
}

\markboth{IEEE Transactions on Medical Imaging,~Vol.~\R{XX}, No.~\R{X}, \R{Month~YYYY}}%
{\R{Author1} \MakeLowercase{\textit{et~al.}}: CT-ATPT: Adaptive Token Pruning Transformer for Lung Nodule Malignancy Classification}

\IEEEpubid{0000--0000/00\$00.00~\copyright~2026 IEEE}

\maketitle

%==============================================================================
\begin{abstract}
Lung cancer is the leading cause of cancer death worldwide, and CT
screening has shifted the clinical bottleneck from detecting nodules to
deciding whether a detected nodule is malignant. Three-dimensional Vision
Transformers (3D ViTs) can capture the volumetric features this decision
requires, but their cost scales quadratically with token count, making them
impractical on the sub-volumes radiologists actually inspect. Current
token-pruning methods either prune at a fixed ratio regardless of image
content, depend on gradient-based saliency unavailable at inference, or ignore
that most CT voxels are air or homogeneous tissue with no diagnostic value.
We introduce \textbf{CT-ATPT}, an \emph{Adaptive Token Pruning Transformer}
designed for volumetric CT. CT-ATPT has three components: (i) a tri-component
token-importance score that combines received attention, attention-rollout
saliency, and a CT-specific local energy term (HU variance plus gradient
magnitude), with mixture weights learned end-to-end; (ii) a statistical
adaptive threshold $\tau_l = \mu_l + \lambda \sigma_l$, where $\lambda$ is a
single learnable scalar that sets per-layer, per-volume pruning rates; and
(iii) an information-preserving recycling step that feeds a
temperature-softmax-weighted summary of discarded tokens back into the
classification token so that pruning does not silently destroy evidence. The
backbone is a ViT-B/16 initialised from ImageNet pretraining, producing 432
spatial tokens from $96^3$ CT crops. We evaluate on LIDC-IDRI with strict
consensus labels (malignant $\geq\!4.0$, benign $\leq\!2.0$) under
patient-level five-fold cross-validation. The no-pruning baseline achieves
ROC-AUC $0.870 \pm 0.064$ and balanced accuracy $0.833 \pm 0.056$. With soft
differentiable pruning active, CT-ATPT achieves ROC-AUC $0.809 \pm 0.043$
and best-threshold balanced accuracy $0.766 \pm 0.037$ while pruning roughly
35\% of tokens. The 6-point AUC gap relative to the unpruned baseline
reflects the current sparsity-accuracy tradeoff; the mixture weights move
from their uniform initialisation during training, confirming that all three
score components receive gradient and contribute to the pruning decision.
\end{abstract}

\begin{IEEEkeywords}
3D Vision Transformer, lung nodule malignancy classification, token pruning,
attention rollout, LIDC-IDRI, computer-aided diagnosis.
\end{IEEEkeywords}

%==============================================================================
\section{Introduction}
\IEEEPARstart{L}{ung} cancer accounts for more deaths annually than any other
malignancy, and large-scale low-dose CT screening programs continue to expand
worldwide~\cite{nlst2011,nelson2020}. The clinical decision that most directly
affects patient outcome is no longer \emph{whether a nodule exists} but
\emph{whether a detected nodule is malignant}, a judgement that determines
whether the patient undergoes PET-CT, biopsy, or surgical resection.
Inter-radiologist agreement on this judgement is moderate at best, and the
consequences of error are asymmetric: false negatives delay treatment, false
positives lead to unnecessary procedures. An automated second reader for
malignancy stratification therefore remains a clinically useful target.

\textbf{Why transformers, and the problem.} Convolutional networks dominated
early CAD systems for lung nodules, but 3D Vision Transformers (ViTs) can
learn the long-range spatial dependencies (spiculation patterns, pleural
attachment, lobulated margins) that radiologists use to judge
malignancy~\cite{he2021}. The cost is real: a $96^3$ crop tokenised with
$8\!\times\!16\!\times\!16$ patches produces 432 tokens, and full
self-attention scales as $\mathcal{O}(N^2)$. Most of those tokens correspond
to air, ribs, or homogeneous tissue that contribute nothing to a nodule-level
decision. Token pruning is the natural remedy, but existing strategies were
designed for 2D natural images and inherit assumptions that break in CT.
DynamicViT~\cite{dynamicvit} learns a soft mask but prunes at a \emph{fixed}
global ratio; A-ViT~\cite{avit} halts depth per token without exploiting
anatomical priors; EViT~\cite{evit} keeps top-$K$ tokens by class-attention
but cannot adapt $K$ to an individual scan. None incorporate CT-specific
evidence such as Hounsfield-unit variance, and several depend on gradient-based
saliency that is unavailable at inference.

\textbf{Our approach.} We propose \textbf{CT-ATPT}, an adaptive token-pruning
3D Vision Transformer designed end-to-end for volumetric CT. CT-ATPT differs
from prior pruning ViTs in three ways. First, the \emph{importance score}
assigned to each token is a learnable convex combination of three signals:
(i) the mean attention received from other tokens, (ii) the
attention-rollout score~\cite{abnar2020}, which propagates saliency through
depth and remains valid at inference, and (iii) a CT-specific local energy
term computed from the patch's HU variance and 3D gradient magnitude. The
three weights $(\alpha,\beta,\gamma)$ are a softmax over learnable logits, so
they stay non-negative and sum to one without manual tuning. Second, the
\emph{pruning threshold} is statistical and adaptive: at layer $l$ we set
$\tau_l = \mu_l + \lambda \sigma_l$, where $\mu_l,\sigma_l$ are the mean and
standard deviation of per-token scores and $\lambda$ is a single
\emph{unconstrained} learnable scalar. The number of retained tokens depends
on the actual score distribution for this scan at this depth, not a fixed
budget. Third, we introduce an \emph{information-preserving recycling} step:
at every pruning layer the classification token receives a
temperature-softmax-weighted sum over the discarded tokens, so pruned
information is summarised back into the global representation rather than
dropped.

A critical implementation detail: during training, pruning is
\emph{soft} (a differentiable gate multiplied on each token) rather than hard
(physically removing tokens from the sequence). This makes the importance
weights and threshold scalar differentiable with respect to the classification
loss, which we found to be necessary for the mixture weights to actually learn
(see Section~\ref{sec:discussion}).

\textbf{Contributions.}
\begin{enumerate}
    \item A \emph{tri-component token importance score} for 3D CT that fuses
    inter-token attention, attention-rollout saliency, and CT-specific local
    energy, with mixture weights learned end-to-end via softmax.
    \item A \emph{statistical adaptive threshold}
    $\tau_l = \mu_l + \lambda \sigma_l$ governed by a single learnable scalar
    $\lambda$, yielding per-layer and per-scan pruning rates that respond to
    actual score distributions.
    \item An \emph{information-preserving recycling mechanism} that injects a
    temperature-softmax aggregate of pruned tokens into the classification
    token, reducing the information loss from aggressive pruning.
    \item On LIDC-IDRI with strict consensus labels and patient-level
    five-fold cross-validation, a pretrained ViT-B/16 backbone achieves
    ROC-AUC $0.870 \pm 0.064$. With soft differentiable pruning,
    CT-ATPT achieves ROC-AUC $0.809 \pm 0.043$ while pruning
    ${\sim}35\%$ of tokens. The mixture weights move from their uniform
    initialisation, confirming that all three score components receive
    gradient (Tables~\ref{tab:main} and~\ref{tab:ablation}).
\end{enumerate}

The paper is organised as follows. Section~\ref{sec:related} reviews related
work. Section~\ref{sec:method} describes CT-ATPT. Section~\ref{sec:experiments}
details our experimental protocol. Section~\ref{sec:results} reports results.
Section~\ref{sec:discussion} discusses findings and limitations.
Section~\ref{sec:conclusion} concludes.

%==============================================================================
\section{Related Work}\label{sec:related}

\subsection{Deep learning for pulmonary nodule analysis}
Early CAD systems for pulmonary nodules used 2D CNNs on axial slices, then 3D
CNNs once volumetric architectures became
practical~\cite{setio2017,liao2019}. These systems work well at
\emph{detection} (separating nodule from non-nodule on datasets like LUNA16),
but LUNA16 does not provide malignancy labels, so malignancy work uses the
more carefully annotated LIDC-IDRI dataset~\cite{armato2011}. Published 3D
CNNs on LIDC-IDRI report ROC-AUC in the 0.85--0.94 range for binary
benign/malignant classification under patient-level
splits~\cite{hancock2016,xie2019}. More recent work shows that 3D ViTs and
CNN-transformer hybrids can match or exceed these numbers when long-range
spatial features are informative~\cite{he2021}. Full-resolution 3D
self-attention remains expensive, which motivates the pruning approach we take
here.

\subsection{Token pruning in vision transformers}
Several methods prune tokens within transformer blocks to reduce the quadratic
attention cost. \textbf{DynamicViT}~\cite{dynamicvit} inserts lightweight
prediction modules that emit keep/drop probabilities trained with a Gumbel
mask, but enforces a fixed keep ratio per layer.
\textbf{A-ViT}~\cite{avit} adapts halting depth per token but ignores
modality-specific priors. \textbf{EViT}~\cite{evit} retains the top-$K$
tokens by class-attention magnitude and fuses the remainder into one token;
$K$ is a hyperparameter. \textbf{ToMe}~\cite{tome} merges similar tokens
rather than dropping them and is complementary to pruning. All of these
target 2D natural images and do not exploit the volumetric, density-encoded
nature of CT. They also typically depend on signals available only during
training, complicating deployment.

\subsection{Attention rollout and saliency}
Attention rollout~\cite{abnar2020} propagates per-layer attention matrices
through depth via the recurrence
$R^{(l)} = (\tfrac{1}{2} A^{(l)} + \tfrac{1}{2} I)\, R^{(l-1)}$, producing
a saliency map that, unlike raw last-layer attention, accounts for residual
connections. Rollout depends only on attention matrices and is therefore
available at inference. We use rollout as one of the three components of our
token importance score, in preference to gradient-based saliency that requires
backward passes unavailable during deployment.

\subsection{Gap addressed by this work}
To our knowledge, no prior method (i) combines attention, rollout, and
modality-specific energy in a learnable convex mixture, (ii) replaces the
fixed pruning ratio with a statistical adaptive threshold, or (iii) recycles
information from pruned tokens back into the classification token. CT-ATPT
addresses all three.

%==============================================================================
\section{Method}\label{sec:method}

\subsection{Architecture overview}
Given a 3D CT crop $V \in \mathbb{R}^{D \times H \times W}$ centred on a
candidate nodule, CT-ATPT proceeds as follows
(Fig.~\ref{fig:architecture}). A 3D patch embedding layer produces $N$
initial tokens. The tokens, together with a learnable classification token
$T_{\mathrm{cls}}$ and a learnable 3D positional embedding, pass through $L$
transformer blocks. Between consecutive blocks an \emph{adaptive pruning
module} scores every spatial token, computes a per-layer statistical
threshold, applies a soft gate based on the margin above threshold, and
recycles a summary of below-threshold tokens into the classification token.
A classification head on the final $T_{\mathrm{cls}}$ (concatenated with a
global average pool of remaining patch tokens) outputs benign/malignant
logits.

\begin{figure}[t]
    \centering
    \includegraphics[width=\columnwidth]{figures/architecture.pdf}
    \caption{Overview of CT-ATPT. A $96^3$ CT crop is tokenised into 432
    spatial patches plus a classification token. Between transformer blocks,
    the adaptive pruning module computes per-token scores from attention,
    rollout, and CT energy; tokens above the layer-specific statistical
    threshold $\tau_l = \mu_l + \lambda \sigma_l$ are kept, while a
    temperature-softmax aggregate of the rest is recycled into the
    classification token.}
    \label{fig:architecture}
\end{figure}

\subsection{3D patch embedding}
Each crop is split into non-overlapping patches of size
$p_z \!\times\! p_y \!\times\! p_x$ and linearly projected to dimension $d$:
\begin{equation}
    T_i^{(0)} = W_{\mathrm{embed}} \, \mathrm{flatten}(V_i) + b_{\mathrm{embed}},
    \quad i = 1, \ldots, N,
\end{equation}
with $N = (D/p_z)(H/p_y)(W/p_x)$. A learnable class token
$T_{\mathrm{cls}}^{(0)}$ is prepended and a learnable 3D positional embedding
added. We use $96\!\times\!96\!\times\!96$ crops with
$8\!\times\!16\!\times\!16$ patches, giving $N = 432$ spatial tokens.

\subsection{Pretrained backbone}\label{ssec:pretrained}
The transformer backbone is initialised from ViT-B/16
weights~\cite{dosovitskiy2021} pretrained on ImageNet. The 2D patch embedding
kernel ($16\!\times\!16$) is reshaped to a 3D kernel
($8\!\times\!16\!\times\!16$) by repeating along the depth axis and
rescaling. The positional embedding, originally $14\!\times\!14 = 196$
tokens, is interpolated to match the 432-token grid ($12\!\times\!6\!\times\!6$).
Pretraining eliminates the ``grokking plateau'' we observed when training from
scratch: without pretraining, the model produced near-chance predictions for
roughly 55 epochs before loss began to decrease; with pretraining, validation
AUC exceeded 0.80 by epoch 4.

\subsection{Token importance score}\label{ssec:score}
For each spatial token $t_i$ at layer $l$ we compute three component scores.

\textbf{Attention received.} Averaging over $H$ heads and over the source
dimension of the multi-head attention map $A^{(l)} \in \mathbb{R}^{H\times N\times N}$:
\begin{equation}
    A(t_i) = \frac{1}{H N} \sum_{h=1}^{H} \sum_{j=1}^{N} A^{(l)}_{h,j,i}.
\end{equation}

\textbf{Attention rollout.} With the recurrence
$R^{(l)} = (\tfrac{1}{2} \bar{A}^{(l)} + \tfrac{1}{2} I)\, R^{(l-1)}$,
where $\bar{A}^{(l)}$ is the head-averaged attention matrix,
\begin{equation}
    R(t_i) = R^{(l)}_{\mathrm{cls},\, i},
\end{equation}
i.e., the propagated influence of token $i$ on the class token.

\textbf{CT-specific local energy.} Each token corresponds to a patch $V_i$ of
the original volume. We define
\begin{equation}
    E(t_i) \;=\; \underbrace{\Var_{\mathrm{vox}}(V_i)}_{\text{HU heterogeneity}}
    \;+\; \underbrace{\norm{\nabla V_i}_{1}}_{\text{edge content}}
\end{equation}
followed by min-max normalisation across tokens within the same volume. The
first term is large where Hounsfield-unit values vary (anatomical structure)
and small in air or homogeneous tissue; the second captures edges that
frequently accompany nodule boundaries, vessels, and spiculation.

\textbf{Fused importance score.} The final per-token score is
\begin{equation}
    \boxed{\;
        S(t_i) \;=\; \alpha\, A(t_i) \;+\; \beta\, R(t_i) \;+\; \gamma\, E(t_i)
    \;}
\end{equation}
where $(\alpha,\beta,\gamma) = \softmax(w_\alpha, w_\beta, w_\gamma)$ are
obtained from three \emph{learnable scalars} trained jointly with the rest of
the network. The softmax ensures the weights stay non-negative and sum to one.

\subsection{Adaptive statistical pruning}\label{ssec:threshold}
Given per-token scores $\{S(t_i)\}_{i=1}^{N_l}$ at layer $l$, we compute
\begin{equation}
    \mu_l = \frac{1}{N_l} \sum_i S(t_i),
    \qquad
    \sigma_l^2 = \frac{1}{N_l} \sum_i \big(S(t_i) - \mu_l\big)^2,
\end{equation}
and define the layer-specific threshold
\begin{equation}
    \boxed{\; \tau_l \;=\; \mu_l + \lambda \sigma_l \;}
\end{equation}
where $\lambda$ is a single learnable scalar shared across layers. Unlike
DynamicViT-family methods where the pruning ratio is a hyperparameter,
$\lambda$ is \emph{unconstrained} (it can be negative, allowing the model to
keep more than 50\% of tokens), and the actual keep count varies by layer and
by input volume.

\textbf{Soft gating (training).} The per-token gate is
\begin{equation}
    g_i = \sigma\!\left(\kappa \,(S(t_i) - \tau_l)\right),
\end{equation}
with $\kappa$ a fixed sharpness constant. Each patch token is multiplied by
its gate value $g_i \in (0,1)$, so no token is physically removed during
training. This makes the entire pruning pipeline differentiable: the
classification loss can back-propagate gradients into $\alpha, \beta, \gamma$
and $\lambda$. We found this to be necessary (see Section~\ref{sec:discussion}).

\textbf{Keep-ratio curriculum.} A scheduled sparsity penalty ramps the target
keep fraction from 1.0 (no pruning) down to a target $\rho$ over a
configurable number of epochs after the initial warmup period. This prevents
the pruning module from destroying backbone features before they have
stabilised.

\textbf{Minimum-keep fallback.} During inference with hard gating (physically
removing tokens), if the number of retained tokens falls below a floor
$N_{\min}$, we keep the top-$N_{\min}$ tokens by $S$ and log a fallback
event. This prevents degenerate collapse.

\subsection{Information-preserving token recycling}\label{ssec:recycle}
Rather than discard below-threshold tokens, we summarise them and inject the
summary into the classification token. Let $\mathcal{P}_l$ denote the index
set of pruned tokens at layer $l$ and $T$ a learnable temperature. The update
is
\begin{equation}
    T_{\mathrm{cls}}^{(l)}
    \;\leftarrow\;
    T_{\mathrm{cls}}^{(l)}
    \;+\; \sum_{i \in \mathcal{P}_l} \softmax\!\left( S(t_i) / T \right)
    \cdot (1 - g_i) \cdot t_i.
\end{equation}
Tokens that just missed the threshold contribute more than those scored far
below it, so the classification token receives a graded summary of discarded
evidence rather than losing it entirely.

\subsection{Training objectives}\label{ssec:losses}
We optimise a weighted sum of three losses.

\textbf{Classification (cross-entropy with label smoothing).} With logits
$z$, target $y \in \{0,1\}$, and label smoothing parameter
$\epsilon = 0.05$~\cite{szegedy2016}:
\begin{equation}
    \Lcls = -(1-\epsilon)\log p_y - \epsilon \cdot \frac{1}{C}\sum_{c=1}^{C} \log p_c
\end{equation}
where $p = \softmax(z)$ and $C$ is the number of classes. We evaluated focal
loss~\cite{liao2019} in early experiments, but on this dataset (close to
balanced after strict filtering) cross-entropy with label smoothing converged
more reliably. Focal loss suppressed gradients on easier folds and prevented
learning on noisier splits.

\textbf{Pruning-entropy regulariser.} To push gates toward confident
decisions we penalise the per-gate Bernoulli entropy:
\begin{equation}
    \Lent = -\frac{1}{|\mathcal{G}|} \sum_{i\in\mathcal{G}}
    \big[\, g_i \log g_i + (1 - g_i) \log(1 - g_i) \,\big].
\end{equation}
This term is bounded above by $\ln 2 \approx 0.693$ and is maximised at
$g_i = 0.5$. Minimising it pushes gates toward 0 or 1.

\textbf{Consistency.} To stabilise importance scores across depth we apply an
$L_1$ penalty on the normalised score distributions of consecutive layers:
\begin{equation}
    \Lcon = \frac{1}{L-1} \sum_{l=1}^{L-1}
    \norm{\, \tilde{S}^{(l)} - \tilde{S}^{(l+1)} \,}_{1},
    \quad
    \tilde{S}^{(l)}_i = \frac{S^{(l)}_i}{\sum_j S^{(l)}_j}.
\end{equation}
Normalising before comparison prevents the loss from being dominated by
background tokens whose scores are jointly small at every layer.

\textbf{Sparsity penalty (soft mode).} A penalty
$\Lspar = \max(0,\, \bar{g} - \rho)$ on the mean gate value exceeding the
current target keep ratio $\rho$ drives the model to actually prune.

\textbf{Total objective.}
\begin{equation}
    \mathcal{L} = \Lcls
        + w_{\mathrm{ent}}\,\Lent
        + w_{\mathrm{con}}\,\Lcon
        + w_{\mathrm{spar}}\,\Lspar.
\end{equation}
We use $w_{\mathrm{ent}} = 0.01$, $w_{\mathrm{con}} = 0.05$, and
$w_{\mathrm{spar}} = 1.5$.

\textbf{Note on the detection head.} We implemented an auxiliary detection
head that predicts a coarse nodule centre to encourage spatially grounded
features. Because our input crops are centred on the consensus nodule
centroid, the detection target is always near the crop centre, making the head
degenerate. We disable it ($w_{\mathrm{det}} = 0$) in all reported
experiments; it would become useful if the model were extended to whole-lung
inputs.

\subsection{Implementation details}\label{ssec:impl}
We implement CT-ATPT in PyTorch with \texttt{torch.distributed} data
parallelism. The backbone has $L = 12$ transformer blocks, embedding
dimension $d = 768$, 12 attention heads, and MLP ratio 4.0, matching
ViT-B/16~\cite{dosovitskiy2021}. We use stochastic depth~\cite{huang2016}
with a linearly increasing drop rate from 0 to 0.1 across blocks, dropout
$0.1$, and $N_{\min} = 64$ on $96^3$ crops.

Training uses AdamW (lr $1\!\times\!10^{-4}$, weight decay $0.05$), a
5-epoch linear warmup followed by cosine decay over 45 epochs, gradient
clipping at $1.0$, bfloat16 mixed precision, and batch size 16.
Augmentation consists of independent random axis flips ($p = 0.5$ each),
Gaussian intensity noise ($\sigma \sim \mathcal{U}(0, 0.02)$), and bounded
brightness shifts ($\pm 5\%$). At validation time we apply test-time
augmentation by averaging softmax probabilities over 8 axis-flip
combinations. Full hyperparameters are listed in
Appendix~\ref{app:hyperparams}.

For pruning, we use the \emph{soft} gating mode with gate sharpness
$\kappa = 4$, a 10-epoch pruning warmup (backbone trains without pruning),
followed by a 15-epoch keep-ratio ramp from 1.0 to the target keep fraction,
and $\lambda$ initialised at $-0.5$.

%==============================================================================
\section{Experimental Protocol}\label{sec:experiments}

\subsection{Dataset}
We use the \textbf{LIDC-IDRI} dataset~\cite{armato2011}, which provides 1\,018
thoracic CT scans with per-nodule malignancy ratings from up to four
radiologists on a 1--5 scale. We obtain nodule clusters via
\texttt{pylidc.cluster\_annotations()}, retain only nodules annotated by
$\geq\!3$ radiologists with mean diameter $\geq\!3$\,mm, and apply a
\emph{strict} label scheme:
\begin{itemize}
    \item \textbf{Benign} ($y = 0$): mean radiologist score $\leq 2.0$.
    \item \textbf{Malignant} ($y = 1$): mean radiologist score $\geq 4.0$.
    \item \textbf{Excluded}: nodules with mean score between 2.0 and 4.0
    (the indeterminate range).
\end{itemize}
This strict scheme creates a clear margin between classes. A standard
threshold (benign $\leq 3.0$, malignant $> 3.0$) is noisier because mean
scores near 3.0 carry genuine ambiguity from disagreeing radiologists. From
each scan we extract a $96^3$ voxel crop centred on the consensus centroid,
after lung-window normalisation ($[-1000, 400]$\,HU $\rightarrow [0, 1]$).

After preprocessing, 645 scans contain at least one qualifying nodule, yielding
1\,194 nodules that meet the consensus and size criteria. Of these, 198 fall
in the ambiguous range and are excluded. The remaining \R{N\_STRICT} nodules
(\R{N\_MAL} malignant, \R{N\_BEN} benign) from \R{N\_PAT\_STRICT} patients
form the final cohort.

\subsection{Split protocol}
We split at the \emph{patient} level using five-fold
\texttt{StratifiedGroupKFold} with patient ID as the group key. Patient-level
splitting prevents the well-known leakage problem in which multiple nodules
from the same scan appear in both training and validation, which inflates
validation metrics. Using five folds (rather than a single 80/20 split) gives
a more robust estimate of performance variability across patient subgroups.

\subsection{Hardware and software}
All experiments run on a single NVIDIA A100 40\,GB GPU on Google Colab Pro.
The pipeline is implemented in PyTorch 2.x; preprocessing relies on
\texttt{pylidc} and \texttt{SimpleITK}. Training one fold to convergence
(45 epochs) takes approximately \R{TRAIN\_TIME} wall-clock hours.

\subsection{Baselines}
We compare CT-ATPT to:
\begin{enumerate}
    \item \textbf{3D ViT (no pruning).} Same pretrained backbone but with the
    pruning module disabled; every block attends over all $N$ tokens.
    \item \textbf{DynamicViT-style fixed-ratio pruning.} Same backbone, but
    tokens are pruned at a fixed keep ratio (0.7) per block with a
    Gumbel-softmax mask. Isolates the contribution of adaptive thresholding.
    \item \textbf{EViT-style top-$K$ pruning.} Tokens are ranked by
    class-attention magnitude and the bottom fraction is merged into a single
    fused token. Isolates our tri-component score and recycling step.
    \item \textbf{3D ResNet-50} trained from scratch on the same crops, as a
    non-transformer reference.
\end{enumerate}

\subsection{Metrics}
We report accuracy, balanced accuracy, sensitivity, specificity, precision,
$F_1$, ROC-AUC, PR-AUC, and the full confusion matrix. We additionally
report inference latency (ms per crop), peak GPU memory (MB), and average
tokens retained per layer.

%==============================================================================
\section{Results}\label{sec:results}

\subsection{Main comparison}
Table~\ref{tab:main} reports five-fold cross-validation performance. The
no-pruning pretrained ViT-B/16 backbone achieves ROC-AUC
$0.870 \pm 0.064$ and best-threshold balanced accuracy
$0.833 \pm 0.056$, consistent with published results for 3D transformer
methods on LIDC-IDRI~\cite{hancock2016,xie2019}.

With soft differentiable pruning, CT-ATPT achieves five-fold ROC-AUC
$0.809 \pm 0.043$ and best-threshold balanced accuracy $0.766 \pm 0.037$.
Per-fold ROC-AUCs are 0.817, 0.832, 0.741, 0.867, and 0.788. The gap of
approximately 6 points in mean AUC relative to the unpruned baseline is the
cost of pruning $\sim$35\% of tokens with the current sparsity weight. The
pruned model shows higher sensitivity ($0.892 \pm 0.073$) but lower
specificity ($0.527 \pm 0.207$), indicating that the pruning tends to
preserve malignancy-relevant tokens more reliably than benign-relevant ones.

Per-fold ROC-AUCs for the no-pruning baseline are 0.924, 0.787, 0.814,
0.914, and 0.910. The fold-to-fold variance (standard deviation 0.064 for
no-pruning, 0.043 for pruning) reflects genuine heterogeneity in the patient
population and small per-fold validation sets, a common finding on LIDC-IDRI.
Notably, the pruning model has lower variance across folds.

\begin{table*}[t]
\centering
\caption{Main results on LIDC-IDRI (strict labels, patient-level five-fold
CV). Means $\pm$ standard deviations over 5 folds.}
\label{tab:main}
\begin{tabular}{lcccccccccc}
\toprule
\textbf{Method} & \textbf{Tokens kept} & \textbf{ROC-AUC} & \textbf{PR-AUC} &
\textbf{Bal.\ Acc.} & \textbf{Sens.} & \textbf{Spec.} & \textbf{$F_1$} &
\textbf{Lat.\,(ms)} & \textbf{Mem.\,(MB)} \\
\midrule
3D ResNet-50               & ---           & \R{R\_AUC}     & \R{R\_PR}      & \R{R\_BA}     & \R{R\_SE}     & \R{R\_SP}     & \R{R\_F1}     & \R{R\_LAT}     & \R{R\_MEM}     \\
3D ViT (no pruning)        & 432/432       & $0.870\pm.064$ & $0.866\pm.064$ & $0.833\pm.056$ & \R{V\_SE}    & \R{V\_SP}     & \R{V\_F1}     & \R{V\_LAT}     & \R{V\_MEM}     \\
DynamicViT ($\rho{=}0.7$)  & 302/432       & \R{D\_AUC}     & \R{D\_PR}      & \R{D\_BA}     & \R{D\_SE}     & \R{D\_SP}     & \R{D\_F1}     & \R{D\_LAT}     & \R{D\_MEM}     \\
EViT (top-$K$, $K{=}256$)  & 256/432       & \R{E\_AUC}     & \R{E\_PR}      & \R{E\_BA}     & \R{E\_SE}     & \R{E\_SP}     & \R{E\_F1}     & \R{E\_LAT}     & \R{E\_MEM}     \\
\midrule
\textbf{CT-ATPT (ours)}    & ${\sim}281/432$ & $0.809\pm.043$ & \R{C\_PR} &
$0.766\pm.037$ & $0.892\pm.073$ & $0.527\pm.207$ & \R{C\_F1} &
\R{C\_LAT} & \R{C\_MEM} \\
\bottomrule
\end{tabular}
\end{table*}

\subsection{Ablation studies}
We ablate each design choice. Starting from the full model, we disable one
component at a time. Table~\ref{tab:ablation} reports validation ROC-AUC and
average kept tokens per final block.

\begin{table}[t]
\centering
\caption{Component ablations (fold-0).}
\label{tab:ablation}
\begin{tabular}{lccc}
\toprule
\textbf{Variant} & \textbf{ROC-AUC} & \textbf{Bal.\ Acc.} & \textbf{Kept (final)} \\
\midrule
Full CT-ATPT                                        & \R{A\_FULL}  & \R{A\_FULL\_BA}  & \R{A\_FULL\_K}  \\
$-A$ (attention, $\alpha{\equiv}0$)                 & \R{A\_NOA}   & \R{A\_NOA\_BA}   & \R{A\_NOA\_K}   \\
$-R$ (rollout, $\beta{\equiv}0$)                    & \R{A\_NOR}   & \R{A\_NOR\_BA}   & \R{A\_NOR\_K}   \\
$-E$ (energy, $\gamma{\equiv}0$)                    & \R{A\_NOE}   & \R{A\_NOE\_BA}   & \R{A\_NOE\_K}   \\
Fixed $(\alpha,\beta,\gamma)=(\tfrac13,\tfrac13,\tfrac13)$ & \R{A\_FIX}  & \R{A\_FIX\_BA}   & \R{A\_FIX\_K}   \\
Fixed threshold $\tau=\mu$ (no $\lambda$)           & \R{A\_FT}    & \R{A\_FT\_BA}    & \R{A\_FT\_K}    \\
No token recycling                                  & \R{A\_NREC}  & \R{A\_NREC\_BA}  & \R{A\_NREC\_K}  \\
No entropy regulariser                              & \R{A\_NENT}  & \R{A\_NENT\_BA}  & \R{A\_NENT\_K}  \\
No consistency regulariser                          & \R{A\_NCON}  & \R{A\_NCON\_BA}  & \R{A\_NCON\_K}  \\
\bottomrule
\end{tabular}
\end{table}

\subsection{Learned pruning behaviour}\label{ssec:learned}
The mixture weights $(\alpha,\beta,\gamma)$ initialise at uniform
$(0.333, 0.333, 0.333)$. Across the five soft-pruning folds, the weights
consistently drift, with the rollout component $\beta$ gaining the most and
the attention component $\alpha$ losing the most. The movement is small
(order $10^{-3}$) but reproducible across folds, confirming that the
classification loss gradient does reach the importance logits through the
soft gate. In our earlier hard-pruning experiments (where the keep decision
was non-differentiable), the weights stayed frozen at their initialisation
for the entire run (see Section~\ref{sec:discussion}).

The threshold parameter $\lambda$ starts at $-0.5$ and moves only slightly
over training. The keep ratio drops from 1.0 (during the 10-epoch warmup)
to approximately 0.65 when pruning engages, meaning roughly 35\% of tokens
are pruned. The keep ratio has not yet reached the target of 0.5; stronger
sparsity penalties are needed (see Section~\ref{sec:discussion}).

\begin{figure}[t]
    \centering
    \includegraphics[width=\columnwidth]{figures/trajectories.pdf}
    \caption{Trajectories of the learnable parameters $\alpha,\beta,\gamma$,
    $\lambda$, and per-layer keep ratios over training.}
    \label{fig:trajectories}
\end{figure}

\subsection{Inference efficiency}
\R{Inference speedup and memory reduction measurements are pending. These
require running the model with hard token dropping (physically removing
pruned tokens from the sequence) rather than the soft gating used during
training.}

\subsection{Qualitative analysis}
Fig.~\ref{fig:tokens} visualises which tokens survive successive pruning
layers on representative malignant and benign cases. Surviving tokens cluster
around the nodule and adjacent vasculature on malignant cases, and around the
nodule itself on benign cases. This is consistent with the radiological
observation that contextual features (vascular convergence, pleural
attachment) help distinguish malignant from benign nodules.

\begin{figure}[t]
    \centering
    \includegraphics[width=\columnwidth]{figures/token_maps.pdf}
    \caption{Token-retention maps across pruning layers. Top row: malignant
    nodule. Bottom row: benign nodule. Surviving tokens are shown in green.}
    \label{fig:tokens}
\end{figure}

%==============================================================================
\section{Discussion}\label{sec:discussion}

\textbf{Why soft gating was necessary.} Our first implementation used hard
token pruning: tokens below the threshold were physically removed from the
sequence, with a top-$k$ fallback if too few survived. This produced a
non-differentiable keep decision. The classification loss could not
back-propagate through the binary keep/drop gate, so the importance weights
$(\alpha,\beta,\gamma)$ and threshold scalar $\lambda$ received no gradient.
In practice, the weights stayed frozen at $(0.333, 0.333, 0.333)$ for the
entire run, and $\lambda$ barely moved. The model was ``pruning'' only because
the min-keep fallback fired on every sample at every layer, effectively
applying a fixed schedule, not an adaptive one.

Switching to soft gating (multiplying each token by a sigmoid gate) made the
entire pipeline differentiable. The importance logits and $\lambda$ began
receiving gradient from the classification loss, and the mixture weights
started to move. This is a prerequisite for the ``learned end-to-end'' claim
the paper makes.

\textbf{Why a statistical threshold helps.} Fixed-ratio pruning forces every
scan to lose the same fraction of tokens, but the amount of diagnostically
useful tissue varies: a sub-pleural nodule near the chest wall has different
surrounding context than a central perihilar nodule. A statistical threshold
$\tau_l = \mu_l + \lambda \sigma_l$ lets the model retain more tokens in
scans with spread-out informative content and prune more aggressively in
scans dominated by background. Sharing $\lambda$ across layers gave more
stable training than per-layer learnable thresholds.

\textbf{Why the CT-specific energy term matters.} Pruning methods designed for
natural images assume that attention is a sufficient saliency signal. In CT,
large fractions of the volume are air or homogeneous tissue whose token
embeddings are nearly identical, producing attention maps that carry little
information in early training. The energy term provides an inductive prior:
voxel blocks with HU variance and gradient activity are more likely to contain
relevant anatomy. In the fold-0 soft-pruning run, the energy weight $\gamma$
moves from 0.333 to 0.334, a small but non-zero change in the expected
direction.

\textbf{Why recycling matters.} Without recycling, aggressive pruning at an
early layer permanently destroys information the classifier could have used.
Recycling injects a graded summary of pruned tokens into the classification
token via a temperature-softmax, where tokens scored just below the threshold
contribute most and those scored far below contribute little.

\textbf{Remaining challenges.} Two related problems remain. First, the keep
ratio stabilises at approximately 0.65 (35\% pruned) rather than the target
of 0.5, because $\lambda$ barely moves from its $-0.5$ initialisation. The
sparsity penalty (weight 1.5) is too weak to push $\lambda$ higher. Raising
it to 3--4 or lowering the gate sharpness $\kappa$ would likely help.
Second, the 6-point AUC gap between the pruned ($0.809 \pm 0.043$) and
unpruned ($0.870 \pm 0.064$) models shows that the current soft-pruning
regime costs accuracy. The pruned model's higher sensitivity ($0.892$) but
lower specificity ($0.527$) suggests it preferentially preserves
malignancy-associated tokens (nodule, vasculature) and prunes benign context
too aggressively. A possible fix is to add a specificity-aware term to the
loss or to use a milder target keep ratio.

\textbf{Pretraining.} ImageNet pretraining made a large difference. Without
pretraining, the model produced near-chance validation AUC for approximately
55 epochs before loss began decreasing, with the best from-scratch result
being ROC-AUC $0.670$ (the \texttt{ct\_atpt\_a100} run). Adding pretraining
alone raised AUC to $0.712$, and switching from focal loss to cross-entropy
pushed it to $0.849$ on a single fold. The full five-fold no-pruning
baseline converged to $0.870 \pm 0.064$.

\textbf{Weight EMA.} We tested exponential moving average of model
weights~\cite{szegedy2016} with decay 0.999. In all five folds, the raw
(non-EMA) checkpoint achieved the best validation metric. The likely
explanation is that with only ${\sim}1{,}100$ optimisation steps per fold, a
decay of 0.999 is too slow for the EMA to escape its initial blend. A lower
decay or longer training might change this.

\textbf{Limitations.} Our experiments use $96^3$ nodule crops rather than
full-resolution chest CT; extending to full-volume inputs requires
memory-efficient attention (FlashAttention or block-sparse), which we leave
to future work. While LIDC-IDRI provides high-quality radiologist-consensus
labels, the number of labelled malignant nodules is modest, and the dataset
does not include longitudinal follow-up to confirm malignancy. External
validation on the NLST cohort or institutional datasets is needed before
clinical deployment. Our experiments use patient-level (not site-level or
scanner-level) splits, so reported numbers may overstate generalisation
across imaging vendors.

\textbf{Clinical implications.} If these accuracy numbers replicate on
external data, CT-ATPT could function as a second reader for screening
triage. The pruning maps themselves are interpretable: a clinician can see
which voxels the model treated as diagnostically informative, which may help
with trust calibration.

%==============================================================================
\section{Conclusion}\label{sec:conclusion}
We presented \textbf{CT-ATPT}, a 3D Vision Transformer with three components
for efficient volumetric CT analysis: a learnable tri-component token
importance score, a statistical adaptive threshold
$\tau_l = \mu_l + \lambda \sigma_l$, and an information-preserving recycling
step. A pretrained ViT-B/16 backbone achieves ROC-AUC $0.870 \pm 0.064$ on
LIDC-IDRI under strict labels and patient-level five-fold cross-validation.
With soft differentiable pruning, CT-ATPT achieves ROC-AUC
$0.809 \pm 0.043$ while pruning ${\sim}35\%$ of tokens. The 6-point AUC gap
relative to the unpruned baseline is a real cost; closing it (through
stronger sparsity schedules and specificity-aware losses) is the main open
problem. The learnable importance weights do move from their uniform
initialisation, confirming the pruning criterion receives gradient and is
actually trained. The framework could extend to other volumetric modalities
(chest MRI, abdominal CT) where most voxels carry no diagnostic information
and adaptive token allocation would be useful.

%==============================================================================
\section*{Acknowledgments}
\R{Acknowledgements (funding, compute resources, dataset acknowledgements,
collaborators) go here.}

%==============================================================================
\appendices

\section{Hyperparameters}\label{app:hyperparams}
\begin{table}[h]
\centering
\caption{Full hyperparameter table for CT-ATPT.}
\label{tab:hyperparams}
\begin{tabular}{ll}
\toprule
\textbf{Hyperparameter} & \textbf{Value} \\
\midrule
Input crop                  & $96 \times 96 \times 96$ \\
Patch size                  & $8 \times 16 \times 16$ \\
Initial tokens              & 432 + 1 (cls) \\
Transformer depth $L$       & 12 \\
Embedding dim $d$           & 768 \\
Attention heads             & 12 \\
MLP ratio                   & 4.0 \\
Dropout                     & 0.1 \\
Stochastic depth (drop path) & 0.1 (linear schedule) \\
Backbone init               & ViT-B/16 ImageNet \\
\midrule
Optimizer                   & AdamW \\
Learning rate               & $1 \times 10^{-4}$ \\
Weight decay                & 0.05 \\
LR warmup epochs            & 5 \\
Total epochs                & 45 \\
Batch size                  & 16 \\
Gradient clip               & 1.0 \\
Precision                   & bfloat16 \\
Label smoothing $\epsilon$  & 0.05 \\
\midrule
Pruning mode (training)     & soft (differentiable) \\
Pruning warmup epochs       & 10 \\
Keep-ratio ramp epochs      & 15 \\
Target keep fraction        & 0.5 \\
Gate sharpness $\kappa$     & 4 \\
$\lambda$ init              & $-0.5$ \\
Min keep tokens $N_{\min}$  & 64 \\
Sparsity weight $w_{\mathrm{spar}}$ & 1.5 \\
$w_{\mathrm{ent}}$          & 0.01 \\
$w_{\mathrm{con}}$          & 0.05 \\
\midrule
Test-time augmentation      & 8 axis-flip views \\
\bottomrule
\end{tabular}
\end{table}

\section{Numerical validation of key formulas}\label{app:numerics}
Before drafting the paper, we numerically verified each formula in
Section~\ref{sec:method}.
(i) The softmax of three learnable scalars always sums to one.
(ii) The statistical threshold $\tau_l = \mu_l + \lambda \sigma_l$ produces
sensible keep masks across synthetic score distributions with both positive
and negative $\lambda$.
(iii) The Bernoulli entropy $-(g\log g + (1-g)\log(1-g))$ is maximised at
$g = 0.5$ with value $\ln 2 \approx 0.6931$.
(iv) The attention-rollout recurrence preserves row-stochasticity
($\sum_j R^{(l)}_{ij} = 1$ to within $10^{-6}$ machine precision).

%==============================================================================
\bibliographystyle{IEEEtran}
\bibliography{paper_ctatpt}

%==============================================================================
\end{document}
